\documentclass[12pt]{article}
\usepackage[utf8]{inputenc}
\usepackage[T1]{fontenc}
\usepackage{hyperref}
\title{Molbog Language Resources for K-12 Curriculum Development}
\author{Compiled by Claude Sonnet 5.5 (claude-sonnet-5-5) for Andr\'as Kornai}
\date{}
\begin{document}
\maketitle

\section{Introduction}
Molbog (ISO 639-3 pwm) is an Austronesian language spoken mainly at the southern tip of Palawan, Philippines (Bataraza and Balabac, with smaller communities in Rizal and Brooke's Point), and in Sabah, Malaysia (Banggi and Balambangan), with some speakers in Mapun, Tawi-Tawi. Wikipedia gives 6,700 speakers (1990 figure); Omniglot gives about 6,000. Almost all speakers are Sunni Muslim. Classification is disputed: Thiessen (1981) and Smith (2017) place it in Palawanic, while Lobel (2013) groups it with Bonggi.

\section{Primary Educational Resources}
No primers, textbooks or school programmes in Molbog were found in search results.

\section{Linguistic and Reference Materials}
\begin{itemize}
\item Wikipedia, ``Molbog language''. \url{https://en.wikipedia.org/wiki/Molbog_language}
\item Wikipedia, ``Molbog-Bonggi languages''. \url{https://en.wikipedia.org/wiki/Molbog-Bonggi_languages}
\item Omniglot page on Molbog. \url{https://mail.omniglot.com/writing/molbog.htm}
\end{itemize}

\section{Religious and Community Resources}
\begin{itemize}
\item Joshua Project, Molbog people group in the Philippines. \url{https://joshuaproject.net/index.php/people_groups/13774/RP}
\end{itemize}

\section{Digital and Audio Resources}
No audio, video or digital corpus resources were seen in search results.

\section{Curriculum Development Context}
Molbog is not known to be taught in schools, and search results show only encyclopedic overviews. A curriculum would require new documentation work with Molbog communities in Palawan and Sabah, and would need to consider the relation to neighbouring languages (Palawanic, Bonggi) and to national and regional languages used in schools.
\end{document}